% \documentclass[12pt,english,ignorenonframetext,aspectratio=169,]{beamer}
\documentclass[11pt,english,ignorenonframetext,]{beamer}

%%%%%%%%%%%%%%%
%% Beamer theme
% choose one from http://deic.uab.es/~iblanes/beamer_gallery/
% or http://www.hartwork.org/beamer-theme-matrix/
% \usetheme{Warsaw}
\usetheme{CambridgeUS}

%%%%%%%%%%%%%%%%%%%%%%
%% Beamer color theme
%% default albatross beaver beetle crane dolphin dove fly lily
%% orchid rose seagull seahorse whale wolverine

%\usecolortheme{seahorse}  %% very lighty
\usecolortheme{dolphin}    %% nice blue
\usecolortheme{orchid}     %% dark red ?
\usecolortheme{whale}      %% black and blue as Warsaw

%%%%%%%%%%%%%%%%%%%%%%%%%%%%%%%%%%%%%%%%%%%%%%%%%%%%%%%%%%%%%%%%%%%%%%%%%%%%%%%%
%% Change the theme
%\setbeamercolor{alerted text}{fg=orange}
%\setbeamercolor{background canvas}{bg=white}
%\setbeamercolor{block body alerted}{bg=normal text.bg!90!black}
%\setbeamercolor{block body}{bg=normal text.bg!90!black}
%\setbeamercolor{block body example}{bg=normal text.bg!90!black}
%\setbeamercolor{block title alerted}{use={normal text,alerted text},fg=alerted text.fg!75!normal text.fg,bg=normal text.bg!75!black}
%\setbeamercolor{block title}{bg=blue}
%\setbeamercolor{block title example}{use={normal text,example text},fg=example text.fg!75!normal text.fg,bg=normal text.bg!75!black}
%\setbeamercolor{fine separation line}{}
\setbeamercolor{frametitle}{fg=black}
%\setbeamercolor{item projected}{fg=black}
%\setbeamercolor{normal text}{bg=black,fg=yellow}
%\setbeamercolor{palette sidebar primary}{use=normal text,fg=normal text.fg}
%\setbeamercolor{palette sidebar quaternary}{use=structure,fg=structure.fg}
%\setbeamercolor{palette sidebar secondary}{use=structure,fg=structure.fg}
%\setbeamercolor{palette sidebar tertiary}{use=normal text,fg=normal text.fg}
%\setbeamercolor{section in sidebar}{fg=brown}
%\setbeamercolor{section in sidebar shaded}{fg= grey}
\setbeamercolor{separation line}{}
%\setbeamercolor{sidebar}{bg=red}
%\setbeamercolor{sidebar}{parent=palette primary}
%\setbeamercolor{structure}{bg=black, fg=green}
%\setbeamercolor{subsection in sidebar}{fg=brown}
%\setbeamercolor{subsection in sidebar shaded}{fg= grey}
%\setbeamercolor{title}{fg=blackblue}
%\setbeamercolor{titlelike}{fg=blackblue}


%%%%%%%%%%%%%%%%%%%%%%%
%% Other beamer options
% \setbeamercovered{transparent}
\setbeamercovered{invisible}
% Permet de laisser en gris le texte qui n'est pas encore apparu (lorsqu'on utilise les commandes avec des <1,2> ou <4-9>.

%\setbeamercolor{normal text}{fg=black,bg=white}

%%%%%%%%%%%%%%%%%%%%%%%
%% Change Beamer fonts
% \usefonttheme{default}
% \usefonttheme[onlymath]{serif}
\usefonttheme{serif}

\setbeamerfont{title}{family=\rm}
\setbeamerfont{titlelike}{family=\rm}
\setbeamerfont{frametitle}{family=\rm}

%%%%%%%%%%%%%%%%%%%%%%%%%%%%%%%%%%%%%%%%%%%%%%%%%%%%%%%%%%%%%%%%%%%%%%%%%%%%%%%%
%% innertheme
%% rectangles circles inmargin rounded
% \useinnertheme{rounded}  % XXX My preference
\useinnertheme{circles}    % XXX

%%%%%%%%%%%%%%%%%%%%%%%%%%%%%%%%%%%%%%%%%%%%%%%%%%%%%%%%%%%%%%%%%%%%%%%%%%%%%%%%
%% outertheme
%% infolines miniframes shadow sidebar smoothbars smoothtree split tree
%\useoutertheme{infolines}

%% No navigation symbol.
\setbeamertemplate{navigation symbols}{}
\beamertemplatenavigationsymbolsempty

% XXX Add a background image to the slides
% \usepackage{tikz}
% \setbeamertemplate{background}{\includegraphics[width=\paperwidth,height=\paperheight,keepaspectratio]{IETR.jpg}}
% \setbeamertemplate{background}{{\centering\begin{tikzpicture}\node[opacity=0.15]{\includegraphics[width=0.98\paperwidth]{IETR_et_partenaires_IETR.png}};\end{tikzpicture}}}

% Other options
%\setbeamertemplate{footline}[page number]

\beamertemplateballitem
\setbeamertemplate{itemize item}[square]


\setbeamertemplate{caption}[numbered]
\setbeamertemplate{caption label separator}{: }
\setbeamercolor{caption name}{fg=normal text.fg}
\beamertemplatenavigationsymbolsempty
\usepackage{lmodern}
\usepackage{color}
  \newcommand{\urlb}[1]{\textcolor{blue}{\url{#1}}}
%% Color definition
\usepackage{xcolor}
%% WARNING attention when changing the colors, change both the {RGB}{r,g,b} and % rgb(r,g,b)
\definecolor{blackblue}{RGB}{19,19,59}  % rgb(48,48,150)
\definecolor{bleu}{RGB}{0,0,204}           % rgb(0,0,204)
\definecolor{deeppurple}{RGB}{102,0,204}   % rgb(102,0,204)
\definecolor{darkgreen}{RGB}{0,100,0}      % rgb(0,100,0)
\definecolor{yellowgreen}{RGB}{200,215,0}  % rgb(200,215,0)
\definecolor{bluegreen}{RGB}{0,185,140}    % rgb(0,185,140)
\definecolor{gold}{RGB}{255,180,0}         % rgb(255,180,0)
\definecolor{meca}{RGB}{255,110,0}         % rgb(255,110,0)
\definecolor{strongred}{RGB}{255,0,0}      % rgb(255,0,0)
\definecolor{normalred}{RGB}{204,0,0}      % rgb(204,0,0)
\definecolor{darkred}{RGB}{174,0,0}        % rgb(174,0,0)
\definecolor{info}{RGB}{174,0,0}        % rgb(174,0,0)
\definecolor{darkblue}{RGB}{0,0,174}        % rgb(0,0,174)
\definecolor{maths}{RGB}{0,0,174}        % rgb(0,0,174)
\definecolor{darkpurple}{RGB}{114,0,114}   % rgb(114,0,114)
\definecolor{ml}{RGB}{114,0,114}   % rgb(114,0,114)

\usepackage{amssymb,amsmath}
\usepackage{bbm,bm}  % bold maths symbols
\usepackage{ifxetex,ifluatex}
\usepackage{fixltx2e} % provides \textsubscript

% % FIXME remove as soon as possible, it slows down compilation to import TikZ
% %% TikZ
% \usepackage{tikz}
% \usetikzlibrary{snakes,arrows,shapes}
% % https://tex.stackexchange.com/a/226974/
% \tikzset{
%   font={\fontsize{10pt}{10}\selectfont}
% }
\usepackage{tikzsymbols}  % https://tex.stackexchange.com/a/227226/97964

\IfFileExists{macrosText.sty}{\usepackage{macrosText}}{}

\usepackage[linesnumbered,commentsnumbered,inoutnumbered,slide]{algorithm2e}


\ifnum 0\ifxetex 1\fi\ifluatex 1\fi=0 % if pdftex
  \usepackage[T1]{fontenc}
  \usepackage[utf8]{inputenc}
\else % if luatex or xelatex
  \ifxetex
    \usepackage{mathspec}
  \else
    \usepackage{fontspec}
  \fi
  \defaultfontfeatures{Ligatures=TeX,Scale=MatchLowercase}
\fi
% use upquote if available, for straight quotes in verbatim environments
\IfFileExists{upquote.sty}{\usepackage{upquote}}{}
% use microtype if available
\IfFileExists{microtype.sty}{%
\usepackage{microtype}
\UseMicrotypeSet[protrusion]{basicmath} % disable protrusion for tt fonts
}{}
\newif\ifbibliography
\hypersetup{
            pdfborder={0 0 0},
            breaklinks=true}
% \urlstyle{same}  % don't use monospace font for urls
% Code embedding.
\usepackage{palatino}              % Use the Palatino font % XXX remove if it is ugly ?
\usepackage{graphicx,grffile}
\makeatletter
\def\maxwidth{\ifdim\Gin@nat@width>\linewidth\linewidth\else\Gin@nat@width\fi}
\def\maxheight{\ifdim\Gin@nat@height>\textheight0.8\textheight\else\Gin@nat@height\fi}
\makeatother
% Scale images if necessary, so that they will not overflow the page
% margins by default, and it is still possible to overwrite the defaults
% using explicit options in \includegraphics[width, height, ...]{}
\setkeys{Gin}{width=\maxwidth,height=\maxheight,keepaspectratio}

\ifxetex
\usepackage{fontspec}
\setmainfont[Ligatures=Historic]{TeX Gyre Pagella}
\newfontfamily\FiraCode{Fira Code}
\setmonofont[Contextuals={Alternate}]{Fira Code}
\newfontfamily\Fontify[Path = ../common/]{Fontify-Regular}
\else
\newcommand{\Fontify}{}
\fi

% Prevent slide breaks in the middle of a paragraph:
\widowpenalties 1 10000
\raggedbottom

% \AtBeginPart{
%   \let\insertpartnumber\relax
%   \let\partname\relax
%   \frame{\partpage}
% }
% \AtBeginSection{
%   \ifbibliography
%   \else
%     \let\insertsectionnumber\relax
%     \let\sectionname\relax
%     \frame{\sectionpage}
%   \fi
% }
% % Si on veut faire apparaître le sommaire courant à chaque nouvelle section.
% \AtBeginSection{
%   \begin{frame}<beamer>%[t]
%   %\frametitle{Local outline}  % Translate the name of the slide
%   \frametitle{\hspace{0pt}}  % Translate the name of the slide
%     \begin{LARGE}  % XXX LARGE font for this slide?
%     \tableofcontents[currentsection, sectionstyle=show/hide, subsectionstyle=hide/hide]
%     \end{LARGE}  % XXX LARGE font for this slide?
%   \end{frame}
% }
% \AtBeginSubsection{
%   \let\insertsubsectionnumber\relax
%   \let\subsectionname\relax
%   \frame{\subsectionpage}
% }

\setlength{\parindent}{0pt}
\setlength{\parskip}{6pt plus 2pt minus 1pt}
\setlength{\emergencystretch}{3em}  % prevent overfull lines
\providecommand{\tightlist}{%
  \setlength{\itemsep}{0pt}\setlength{\parskip}{0pt}}
\setcounter{secnumdepth}{5}

% https://tex.stackexchange.com/a/70495/
\usepackage{appendixnumberbeamer}


% For \justifying command, see https://tex.stackexchange.com/a/148696/
\usepackage{ragged2e}
\addtobeamertemplate{frame begin}{}{\justifying}
\addtobeamertemplate{block begin}{}{\justifying}
\addtobeamertemplate{block alerted begin}{}{\justifying}
\addtobeamertemplate{block example begin}{}{\justifying}
\addtobeamertemplate{itemize body begin}{}{\justifying}
\addtobeamertemplate{itemize item}{}{\justifying}
\addtobeamertemplate{itemize subitem}{}{\justifying}
\addtobeamertemplate{itemize subsubitem}{}{\justifying}
\addtobeamertemplate{enumerate body begin}{}{\justifying}
\addtobeamertemplate{enumerate item}{}{\justifying}
\addtobeamertemplate{enumerate subitem}{}{\justifying}
\addtobeamertemplate{enumerate subsubitem}{}{\justifying}
\addtobeamertemplate{description body begin}{}{\justifying}
\addtobeamertemplate{description item}{}{\justifying}


\title[Audition AGPR à l'ENS Rennes]{Audition : poste d'AGPR à l'ENS de Rennes}
% \subtitle{Mercredi 29 mai 2019}

\author[Lilian Besson]{\Large \textbf{Lilian Besson}}

\institute[]{{\large
  Doctorant et agrégé}{\newline
  \normalsize
  \newline Équipe SCEE, laboratoire IETR, CentraleSupélec à Rennes
  \newline \& Équipe SequeL, laboratoire CRIStAL, Inria à Lille}}

\date{Mercredi $29$ mai $2019$}

\begin{document}
\justifying

\begin{frame}[plain]
\titlepage

% XXX manual inclusion of logos
\begin{center}
  \includegraphics[height=0.18\textheight]{../common/LogoIETR.png}
  \includegraphics[height=0.21\textheight]{../common/LogoCS.png}
  \includegraphics[height=0.18\textheight]{../common/LogoInria.jpg}
\end{center}

\end{frame}


% \section{\hfill{}Plan\hfill{}}

% \begin{frame}{Organisation de l'exposé \hfill{}(20 minutes)}

% \begin{enumerate}
% \item
%   Présentation du candidat
% \vspace*{25pt}

% \item
%   Mes expériences d'enseignement\\
%   Mon projet pédagogique à l'ENS de Rennes
% \vspace*{25pt}

% \item
%   Mes expériences de recherche\\
%   Mon projet de recherche à l'IRISA
% \end{enumerate}

% \end{frame}


\section{\hfill{}1. Présentation du candidat\hfill{}}

\subsection{\hfill{}Plan\hfill{}}

\begin{frame}{\textcolor{gray}{Organisation de l'exposé}}

\begin{enumerate}
\item
\alert{\textbf{%
  Présentation du candidat
}}
\vspace*{25pt}

\item
\textcolor{gray}{
  Mes expériences d'enseignement\\
  Mon projet pédagogique à l'ENS de Rennes
}
\vspace*{25pt}

\item
\textcolor{gray}{
  Mes expériences de recherche\\
  Mon projet de recherche à l'IRISA
}
\end{enumerate}

\end{frame}


% \begin{frame}{1. Présentation du candidat}

% \begin{itemize}
% \item
%   Cursus universitaire
%   % \\
%   % 10 ans d'informatique et de mathématiques
% % \vspace*{20pt}
% % \item
% %   Mes expériences d'enseignements
%   % \\
%   % 650 heures depuis 2013
% \vspace*{20pt}
% \item
%   Expériences en programmation \& développement
% % \vspace*{20pt}
% % \item
% %   Mes expériences de recherche
% \end{itemize}

% % \begin{quote}
% % Cf. mon CV ou mon site web pour plus de détails. \\
% % \(\hookrightarrow\) \href{https://perso.crans.org/besson/}{perso.crans.org/besson}
% % \end{quote}

% \end{frame}


\subsection{\hfill{}1.1. Cursus\hfill{}}

\begin{frame}{Cursus universitaire :\\
  10 ans d'\textcolor{info}{informatique} et de \textcolor{maths}{mathématiques}}

\pause

\begin{itemize}
\tightlist
\item
  2009-11 : classe prépa. \textcolor{maths}{M}PSI \& \textcolor{maths}{M}P\(^*\), au Lycée Thiers, Marseille

\vspace*{10pt}
\pause
\item
  2011-16 : \textbf{double cursus \textcolor{maths}{maths} \textcolor{info}{informatique}}, à l'ENS de Cachan

  \begin{itemize}[<+->]
  \tightlist
  \item
    2011-12 : double licence, \textcolor{info}{informatique} \& \textcolor{maths}{mathématiques}
  \item
    2012-13 : M1 en \textcolor{info}{info. théorique} (MPRI) \& M1 en \textcolor{maths}{maths appliquées}
  \item
    2013-14 : \textcolor{maths}{agrégation maths}, \textcolor{info}{option informatique} ($23^{\text{ème}}$)
  \item
    \textcolor{gray}{2014-15 : césure}
  \item
    2015-16 : M2 en \textcolor{ml}{apprentissage statistique} (MVA, major $1^{\text{er}}/115$)
  \end{itemize}

\pause
\vspace*{10pt}
  \item
  2016-19 : \textbf{doctorat en \textcolor{info}{informatique}}, à Rennes

  \begin{itemize}
  \tightlist
  \item
  % ``\textcolor{ml}{Apprentissage séquentiel} multi-joueurs et bandits multi-bras
  % appliqués aux réseaux de l'Internet des Objets''
    ``Applications d'algorithmes de \textcolor{ml}{bandits multi-bras (MAB)} simple-
    et multi-joueurs pour les réseaux de l'Internet des Objets (IoT)''
  \end{itemize}
\end{itemize}

\end{frame}


\subsection{\hfill{}1.2. Programmation \& développement\hfill{}}

\begin{frame}{Expériences en programmation \& développement}

\begin{itemize}
\item
  Passionné de logiciels libres et de libre diffusion !

  \begin{itemize}\tightlist
    \item
  % $\hookrightarrow$
      Tout ce que je rédige depuis 2012 (pour mes cours, ma recherche etc)\\
      est distribué gratuitement en ligne, et sous licence libre
      % (GPL, MIT)
  \item[]
    \(\hookrightarrow\)
    site personnel
    \href{https://perso.crans.org/besson}{\textcolor{blue}{\texttt{perso.crans.org/besson}}}
    % \&
    % \href{https://naereen.GitHub.io}{\textcolor{blue}{\texttt{naereen.GitHub.io}}}
  \end{itemize}

\pause
\vspace*{7pt}
\item
  \alert{Expert en Python et OCaml}

  \begin{itemize}
  \tightlist
  \item
    Utilisation quotidienne de Python et Jupyter pour mes recherches
  \item
    Habitué à donner des TP en Python ou Ocaml, au niveau L1 à M2
  \end{itemize}

\pause
\vspace*{7pt}
\item
  Mais aussi \LaTeX, Bash, git, Julia\ldots{}

  \begin{itemize}
  \tightlist
  \item
    Actif sur Bitbucket et sur GitHub depuis 2012 \\
    \(\hookrightarrow\)
    \href{https://Bitbucket.org/lbesson}{\textcolor{blue}{\texttt{Bitbucket.org/lbesson}}}
    \&
    \href{https://GitHub.com/Naereen}{\textcolor{blue}{\texttt{GitHub.com/Naereen}}}
  \item
    \(10\) formations techniques ``GouTP'' données durant ma thèse
  % \item
  %   Des projets variés, tous diffusés en ligne\\
  %   \(\hookrightarrow\)
  %   \href{https://perso.Crans.org/besson}{\textcolor{blue}{\texttt{perso.Crans.org/besson}}}
  %   \&
  %   \href{https://naereen.GitHub.io}{\textcolor{blue}{\texttt{naereen.GitHub.io}}}
  \end{itemize}

% \vspace*{10pt}
% \item
%   Capable de créer/maintenir des pages et site web, SQL etc
% \item
%   Maintenance de machines (Windows, GNU Linux) pour ma thèse
\end{itemize}

\end{frame}

% \begin{frame}{%
% \protect\hypertarget{rightarrow-github.comnaereen}{%
% \(\rightarrow\)
% \href{https://GitHub.com/Naereen}{\texttt{GitHub.com/Naereen}}}}

% % \begin{figure}
% \centering
% \includegraphics{figures/apercu_GitHub.png}
% % \caption{bg original 90\%}
% % \end{figure}

% \end{frame}


\section{\hfill{}2. Mon projet pédagogique\hfill{}}

\subsection{\hfill{}Plan\hfill{}}

\begin{frame}{\textcolor{gray}{Organisation de l'exposé}}

\begin{enumerate}
  \item
\textcolor{gray}{
  Présentation du candidat
}
\vspace*{25pt}

% \item
% Aperçu de ma thèse, que je soutiendrai en octobre 2019
% \vspace*{25pt}

\item
\alert{\textbf{%
  Mes expériences d'enseignement\\
  Mon projet pédagogique à l'ENS de Rennes
}}
\vspace*{25pt}

\item
\textcolor{gray}{
  Mes expériences de recherche\\
  Mon projet de recherche à l'IRISA
}
\end{enumerate}

\end{frame}


\subsection{\hfill{}2.1. Enseignements\hfill{}}

\begin{frame}{Mes expériences d'enseignements}

\begin{itemize}
\item
Environ $400$ heures avant ma thèse
\vspace*{20pt}
\item
\(64 \) heures par an pendant ma thèse
\vspace*{20pt}
\item
\textcolor{info}{En informatique} et \textcolor{maths}{en mathématiques}
\vspace*{20pt}
\item
Du lycée aux niveaux master et agrégation
\end{itemize}

\end{frame}


\begin{frame}{Enseignements : environ $400$ heures avant ma thèse}

\begin{itemize}
\tightlist
\item
  2013-14 : cours particuliers \& correction de copies, \textcolor{maths}{en maths} au lycée, à Paris
\vspace*{10pt}
\pause
\item
  2014-15 : participation à la création de \emph{Mahindra Ecole Centrale}

  \begin{itemize}
  \tightlist
  \item
    Une école d'ingénieur-e-s à Hyderabad (Inde)
  \item
    Partenariat franco-indien : Mahindra group \& École Centrale Paris
  \item
    \textcolor{maths}{150h de cours et TD en maths, niveau post-bac (MPSI)}
  \item
  \textcolor{info}{\textbf{En charge de l'introduction à l'informatique (CS101) avec Python}}
  \item
    Conceptions des sujets et corrections d'examens pour $230$ élèves
  \end{itemize}

\vspace*{10pt}
\pause
\item
  2015-16 : \textcolor{maths}{khôlles en maths (BCPST et MP)}, et \textcolor{info}{TD/TP/examens en \textbf{informatique avec Python (MP)}}, Lycée Lakanal
\end{itemize}

\end{frame}

\begin{frame}[plain]{Corrections de TP/TP/DS \textcolor{info}{en Python pour ``l'Informatique pour Tous''} en CPGE (MP)}

% \begin{figure}
\centering
\includegraphics{figures/apercu_Lakanal.png}
% \caption{bg original 100\%}
% \end{figure}

\end{frame}


\begin{frame}{Enseignements : \(64 \) heures par an pendant ma thèse}

\begin{itemize}[<+->]
% \tightlist
\item
  2016-18 : TD/TP/examens en \textcolor{info}{informatique fondamentale}, pour un cours de \textbf{calculabilité et complexité}, à l'ENSAI
\vspace*{15pt}
\item
  Août 2017 \& 2018 : évaluation de rapports de stage des élèves de L3 et M1 du \textcolor{info}{département informatique} de l'ENS Rennes
\vspace*{15pt}
\item
  2016-19 : TP, leçons, oraux blancs, pour la
  préparation à \textbf{\textcolor{info}{l'option informatique} de \textcolor{maths}{l'agrégation de mathématiques}}, à l'ENS Rennes
\end{itemize}

\end{frame}

\begin{frame}[plain]
  {Simulateur de Machine de Turing interactif pour des TP d'\textcolor{info}{informatique fondamentale à l'ENSAI}}

% \begin{figure}
\centering
\includegraphics{figures/apercu_ENSAI.png}
% \caption{bg original 85\%}
% \end{figure}

\end{frame}

\begin{frame}[plain]{Ressources publiées pour la préparation à l'agrégation}

% \begin{figure}
\centering
\includegraphics{figures/apercu_ENS_agreg_0.png}
% \caption{bg original 90\%}
% \end{figure}

\end{frame}

\begin{frame}[plain]{Corrections de \emph{tous} les textes d'annales de modélisation (\textcolor{info}{option info}) en OCaml et Python}

% \begin{figure}
\centering
\includegraphics{figures/apercu_ENS_agreg_1.png}
% \caption{bg original 90\%}
% \end{figure}

\end{frame}

\begin{frame}[plain]{Corrections avec des \emph{notebooks Jupyter} OCaml /Python}

% \begin{figure}
\centering
\includegraphics{figures/apercu_ENS_agreg_3.png}
% \caption{bg original 90\%}
% \end{figure}

Format interactif qui mélange texte, code (entrée/sortie) et figures.

\end{frame}


\subsection{}

\begin{frame}{Mon projet pédagogique à l'ENS de Rennes}

Avec ce poste d'Agrégé Préparateur à l'ENS de Rennes, je souhaite\ldots{}

\pause

\vspace*{20pt}
\begin{itemize}
\tightlist

\item
Rester \textbf{très impliqué dans la préparation à \textcolor{info}{l'option informatique}} de \textcolor{maths}{l'agrégation de mathématiques}
% (``\emph{prépa agrég}'')
\vspace*{15pt}

\pause
\item
\textbf{M'investir dans d'autres formations !}
\vspace*{15pt}

\item
Et participer à la \textbf{création de nouvelles formations}

\end{itemize}

\end{frame}


% \subsection{\hfill{}2.2. Rester impliqup\hfill{}}


\subsection{\hfill{}2.2. Option informatique à l'agrégation de mathématiques\hfill{}}

\begin{frame}{Préparation à \textcolor{info}{l'option informatique} à l'agrégation}

Je maîtrise l'\textbf{intégralité du programme}, et je connais bien le fonctionnement des \textcolor{info}{deux épreuves orales d'informatique}.

\pause
Je peux\ldots{}

\begin{itemize}
\tightlist
\item
  assurer l'entraînement aux oraux (leçons \& oral de modélisation)
\item
  notamment des TP de programmation en OCaml \& Python
% \item
%   garder \(\simeq\) 15 séances d'entraînements pour les épreuves orales
%   de leçon \& modélisation
% \item
%   être jury pour les oraux blancs
  %  (je peux participer pour les leçons aussi)
\item
  donner des rappels/compléments de cours en algorithmique, fondements etc (selon les besoins)
\end{itemize}

\pause
\vspace*{10pt}

\begin{itemize}
\tightlist
\item
  \ldots{} \textbf{prendre en charge la gestion de la formation ?} (si besoin)
  % \begin{quote}
  %   Je pense connaître assez bien les rouages de son
  % organisation, et l'expérience sera très enrichissante.
  % \end{quote}
\end{itemize}

\end{frame}



% \begin{frame}{Besoin de corriger des lacunes ?}

% % Nous pouvons essayer de mieux préparer pour ces points du programme, les connaissances des élèves sont souvent limitées\ldots{}
% Il y a deux semaines, nous avons constaté que les élèves ont de fortes lacunes sur ces points du programme :

% % \begin{figure}
% \centering
% \includegraphics[width=0.95\textwidth]{figures/programme_specifique_option_D_modelisation.png}
% % \caption{100\%}
% % \end{figure}

% \pause
% \vspace*{-10pt}
% \begin{itemize}
%   \tightlist
%   \item
%   Exemple : 2 élèves ne connaissaient pas la résolution ($4$:$3$) ni la taille (ex. $1024\times768$) de l'écran utilisé pendant les oraux blancs !
%   \begin{small}\begin{quote}
%     ``Je dirais de l'ordre d'un million de pixels sur l'axe horizontal\ldots{}?''
%   \end{quote}\end{small}
% \end{itemize}

% \end{frame}



% \begin{frame}{Deux idées de nouveautés}

% \begin{enumerate}
% [1.]
% \item
%   Des rappels de ``culture générale informatique''

%   \begin{itemize}
%   \tightlist
%   \item
%     connaissances de base en architectures logicielle, matérielle etc
%   \item
%     évalué par un QCM plein de petites connaissances
%   % \item
%     % je peux rédiger un long peut-être rédigé avec les élèves de cette année ?
%   \end{itemize}
% \vspace*{20pt}
% \pause
% \item
%   Rappels sur la représentation des nombres et
%   l'arithmétique ``côté machine'' : nombres entiers binaires, flottants (norme IEEE 754) etc

%   \begin{itemize}
%   \tightlist
%   \item
%     beaucoup de textes de l'oral de modélisation traitent de ce sujet
%   \item
%     les élèves doivent connaître les bases
%   \item
%     et savoir expliquer des erreurs classiques telles que\\
%     \begin{center}
%       \texttt{max\_int + 1 < 0} (en OCaml)\\
%       \texttt{0.1 + 0.1 + 0.1 != 0.3} (en Python, OCaml etc)
%     \end{center}
%   % \item
%   %   $\hookrightarrow$ un TP utilisant à la fois Python et OCaml
%   \end{itemize}
%   % \begin{quote}
%   %   Je veux aussi terminer la
%   %   rédaction d'un nouveau texte de modélisation sur les ``posits'', une
%   %   alternative aux flottants IEEE
%   % \end{quote}
% % \item
% %   Raviver l'utilisation du wiki ``Minerve'' pour l'option info\\
% %   \href{https://minerve.ens-rennes.fr/index.php/Le\%C3\%A7ons_d\%27informatique}{\small \texttt{minerve.ens-rennes.fr/index.php/Leçons\_d'informatique}}
% \end{enumerate}

% \end{frame}


\subsection{\hfill{}2.3. D'autres formations à l'ENS Rennes\hfill{}}

\begin{frame}{M'investir dans d'autres formations}

\pause

% \begin{itemize}
% \item
%   Évaluer les rapports \& soutenances de stages de recherche (L3, M1)
%   % m'a beaucoup plu en août 2017 et 2018\ldots{}
%   \\\(\hookrightarrow\)
%   Utile de connaître les expériences ``hors de l'école'' des élèves !

% \pause
% \vspace*{15pt}
% \item
  Je suis très intéressé par d'autres formations :
  % n'importe quel cours d'informatique ou maths, en L3 et M1.

  \begin{itemize}
  % \tightlist
  % \item
  %   Travaux dirigés ou pratiques
    % (ex : analyse numérique avec MATLAB, Python ou Julia, informatique fondamentale avec OCaml, etc)
  \item
    \textcolor{info}{En informatique fondamentale ou appliquée},
    ou \textcolor{maths}{en mathématiques}
  \vspace*{15pt}
  \item
    En L3 ou en M1, \textcolor{info}{au département d'informatique}
  \vspace*{15pt}
  \item
    Mais aussi aux départements \textcolor{maths}{mathématiques} ou \textcolor{meca}{mécatronique} ?
    % (si besoin)
  \vspace*{15pt}
  \item
    \emph{Thèmes} : algorithmique, programmation, fondements informatique,
    optimisation numérique, proba/stats, \textcolor{ml}{\textbf{apprentissage statistique}}\ldots
  \end{itemize}
% \end{itemize}

\end{frame}


\subsection{\hfill{}2.4. Nouvelles formations\hfill{}}


\begin{frame}{Réfléchir à de nouvelles formations ? \hfill{} \only<2->{(1/2)}}

\pause

J'aimerai créer un cours d'\textcolor{ml}{\textbf{apprentissage statistique}} à l'ENS Rennes

\begin{enumerate}
[1.]
\tightlist
\item
  En L3 : cours d'introduction théorique + TP, aux départements \textcolor{info}{informatique} \& \textcolor{maths}{maths} (\& \textcolor{meca}{mécatronique} ?)
  \hfill{}
  \emph{(semestre \#2)}
\item
  En M1 : compléments théoriques + un ``zoo'' des méthodes existantes + un projet de développement, \textcolor{info}{au dpt info}
  \hfill{}
  \emph{(semestre \#1)}
\end{enumerate}

\pause

\begin{exampleblock}{Quelques idées ?}

\begin{itemize}
\tightlist
\item
  TP = un projet en plusieurs morceaux, sur de ``vraies'' données\\
  Ex : sur la ville de Rennes, \href{https://data.rennesmetropole.fr/}{\texttt{data.RennesMetropole.fr}}
\item
  Dès la seconde édition du cours, j'aimerai permettre aux M1
  d'intervenir auprès des L3
  (entre-aide entre les deux promotions)
  % (ex : les faire évaluer un TP ?)
% \item
%   le lier à un groupe de lecture à l'IRISA ?
\end{itemize}

\end{exampleblock}

\end{frame}


\begin{frame}{Réfléchir à de nouvelles formations ? \hfill{} (2/2)}

\begin{itemize}
\item
  Le \textcolor{info}{CAPES d'informatique} sera ouvert dès la rentrée 2019 !
  % \begin{quote}
  %   L'agrégation d'informatique suivra peut-être (un jour)
  % \end{quote}
\item
  L'ISTIC ouvre une formation au CAPES (\textcolor{info}{M1 informatique})
  % conçue en accord avec l'ESPE de Bretagne
\end{itemize}

\vspace*{30pt}

\begin{alertblock}{Et pour l'agrégation d'informatique\ldots{}?}
  \begin{itemize}
    \item
    Pas encore annoncée, mais elle devrait être créée bientôt !
    \pause
    \item
    Je vais garder un œil sur cette formation au CAPES\ldots{}
    \item
    $\implies$ Pour participer à la création de la \textcolor{info}{\textbf{préparation à l'agrégation d'informatique}} à l'ENS de Rennes ?
  \end{itemize}
  \end{alertblock}

% \begin{exampleblock}{\emph{Idée / espoir}}

% Entraide bi-directionelle : agrégation \(\Longleftrightarrow\) CAPES ?

% % \begin{itemize}
% % \tightlist
% % \item
% %   Une leçon préparée par un-e candidat-e à l'agrég sera un
% %   bon (complément de) cours pour les candidats au CAPES
% % \end{itemize}

% % \begin{quote}
% % Le programme du CAPES n'est pas encore disponible\ldots{}
% % \end{quote}

% \end{exampleblock}

\end{frame}



% \begin{frame}{3.3. Créer de nouvelles formations ? (2/2)}

% \begin{block}{« Pré-prépa agrég »}

% \begin{itemize}
% \tightlist
% \item
%   Entre janvier et avril 2019
% \item
%   Avec trois jeunes agrégé-e-s chercheur-euse-s à l'IRISA
%   % et 2 anciens élèves de la prépa agrég
% \item
%   Nous avons donné quelques séances de compléments de cours en maths aux
%   élèves de M1 du département informatique (\(\simeq\) 20h)
% \item
%   6 élèves étaient déjà intéressés pour préparer l'agrég l'an prochain
% \item
%   Rappels de cours, renforcements théoriques, méthodologie (développement au
%   tableau) etc
% \end{itemize}

% \end{block}

% \(\Longrightarrow\) j'aimerais refaire cette formation en 2020, la
% prendre en charge et la renforcer, s'il y a assez d'élèves !

% \begin{itemize}
% \tightlist
% \item
%   Ici aussi : entreaide prépa agrég \(\Longleftrightarrow\) pré-prépa
%   agrég ?
% \end{itemize}

% \end{frame}



\section{\hfill{}3. Mon projet de recherche\hfill{}}

\subsection{\hfill{}Plan\hfill{}}

\begin{frame}{\textcolor{gray}{Organisation de l'exposé}}

\begin{enumerate}
  \item
\textcolor{gray}{
  Présentation du candidat
}
\vspace*{25pt}

% \item
% Aperçu de ma thèse, que je soutiendrai en octobre 2019
% \vspace*{25pt}

\item
\textcolor{gray}{
  Mes expériences d'enseignement\\
  Mon projet pédagogique à l'ENS de Rennes
}
\vspace*{25pt}

\item
\alert{\textbf{%
  Mes expériences de recherche\\
  Mon projet de recherche à l'IRISA
}}
\end{enumerate}

\end{frame}


\subsection{\hfill{}3.1. Mes expériences de recherche\hfill{}}


\begin{frame}{Mes expériences de recherche}

\pause

\begin{itemize}
\item
  2012, 2013, 2016 : à l'ENS de Cachan, trois stages de recherche
\item
  2013, 2014-15, 2016 : un an et demi à l'étranger (R-U, Inde, Suisse)

\pause
\item
  2016-19 : thèse à la \textbf{frontière entre \textcolor{info}{informatique} et \textcolor{maths}{maths}},
  télécommunications et électronique. Co-encadrement par :

  \begin{itemize}
  \tightlist
  \item
    Christophe Moy (Univ Rennes 1, IETR) expert en \emph{radio intelligente}
  \item
    Émilie Kaufmann (CNRS, labo CRIStAL \& Inria Lille), experte en \textcolor{maths}{statistiques, probabilités} et surtout \textcolor{ml}{les \emph{bandits multi-bras}}
  \end{itemize}

  \pause
  Quelques contributions :

  \begin{itemize}
  \tightlist
  \item
    4 articles dans des conférences internationales (2017, 2018, 2019)
  \item
    2 workshops (2019), 1 démo. à une conf. internationale (2018)
  \item
    3 articles en cours de seconde soumission (2 pour des journaux)
  \item
    SMPyBandits, une bibliothèque de simulations de problèmes de bandits multi-bras en Python
    $\hookrightarrow$ \href{https://SMPyBandits.rtfd.io/}{\textcolor{blue}{\texttt{SMPyBandits.rtfd.io}}}
  \end{itemize}
\end{itemize}

\end{frame}


% \subsection{\hfill{}Démonstration d'un problème de bandit\hfill{}}

\begin{frame}{Bandits multi-bras}

  Prise de décisions séquentielles face à un environment incertain

  \centering
  \includegraphics[height=0.65\textheight]{figures/example_of_a_5_arm_bandit_problem.png}

  {\tiny
      $\hookrightarrow$ \href{https://perso.crans.org/besson/phd/MAB_interactive_demo/}{\textcolor{blue}{\texttt{perso.crans.org/besson/phd/MAB\_interactive\_demo/}}}
  }

\end{frame}

% \begin{frame}{%
% \protect\hypertarget{rightarrow-cv.archives-ouvertes.frlilian-besson}{%
% \(\rightarrow\)
% \href{https://cv.archives-ouvertes.fr/lilian-besson}{\texttt{CV.Archives-Ouvertes.fr/Lilian-Besson}}}}

% % \begin{figure}
% \centering
% \includegraphics{figures/apercu_CV_HAL.png}
% % \caption{bg original 90\%}
% % \end{figure}

% \end{frame}


% \begin{frame}{Aperçu de ma thèse - FIXME schéma à dessiner ?}

% \begin{center}
%   ``\alert<4>{\textbf<4>{Apprentissage séquentiel multi-joueurs}} et \alert<3>{\textbf<3>{bandits multi-bras}} appliqués aux réseaux de l'\alert<2>{\textbf<2>{Internet des Objets}}''
% \end{center}

% \begin{exampleblock}<2->{Objets communicants (Internet of Things)}
%   % Communications sans fil, entre des
%   Appareils autonomes qui communiquent avec une station de base lointaine :
%   faible consommation énergétique, faible capacité de calcul etc.
% \end{exampleblock}

% \begin{exampleblock}<3->{Bandits multi-bras pour la radio intelligente}
%   Idée : laisser les objets apprendre d'eux même une stratégie pour accéder efficacement au réseau, via l'apprentissage par reinforcement.
% \end{exampleblock}

% \begin{exampleblock}<4->{Apprentissage séquentiel multi-joueurs décentralisé}
%   Les objets apprennent à collaborer efficacement, sans supervision, et sans communication entre eux.
% \end{exampleblock}

% \end{frame}


\subsection{\hfill{}3.2. Aperçu de ma thèse\hfill{}}


\begin{frame}{Aperçu de ma thèse}

\begin{center}
  % ``Apprentissage séquentiel multi-joueurs et bandits multi-bras appliqués aux réseaux de l'Internet des Objets''
  ``Applications d'algorithmes de bandits multi-bras (MAB) simple-
    et multi-joueurs pour les réseaux de l'Internet des Objets (IoT)''
\end{center}

Thèmes :

\begin{enumerate}
  \item
  Réseaux sans fil
  \begin{itemize}
    \item
    Objets communicants (``Internet of Things'')
    \item
    Radio intelligente
  \end{itemize}
  \item
  Apprentissage statistique par renforcement
  \begin{itemize}
    \item
    Apprentissage séquentiel, bandits multi-bras
    \item
    Apprentissage multi-joueurs décentralisé
    \item
    \alert{Problèmes non stationnaires, stationnaires par morceaux}
  \end{itemize}
\end{enumerate}

\end{frame}


% \begin{frame}{Quelques publications}
%   \setbeamercovered{transparent}
  
%   \begin{enumerate}
%     \uncover<1>{
%     \item
%     Radio intelligente et objets communicants
%     \begin{itemize}
%     \item
%       \href{https://hal.inria.fr/hal-XXX/document}{\textcolor{darkgreen}{``Decentralized Spectrum Learning for IoT Wireless Networks Collision Mitigation''}}, \href{https://moychris.wordpress.com/}{C. Moy} \& \textbf{L. Besson},
%       \href{https://sites.google.com/view/ISIoT2019/}{workshop ISIoT}, 2019.
%     \item
%       \href{https://hal.inria.fr/hal-01575419/document}``{\textcolor{darkgreen}{Multi-Armed Bandit
%       Learning in IoT Networks and non-stationary settings''}},
%       \href{https://remibonnefoi.wordpress.com/}{R. Bonnefoi}, \textbf{L.
%       Besson}, \href{https://moychris.wordpress.com/}{C. Moy},
%       \href{http://chercheurs.lille.inria.fr/ekaufman/research.html}{E.
%       Kaufmann} \& J. Palicot. \href{http://crowncom.org/2017/}{CrownCom},
%       % (\emph{Conference on Cognitive Radio Oriented Wireless Networks}),
%       % Lisbonne, Portugal,
%       septembre 2017,
%       % \href{https://hal.inria.fr/hal-01575419}{HAL-01575419},
%       \alert{\textbf{prix du meilleur article !}}
%     \end{itemize}
%     }

%     \vspace*{10pt}
%     \uncover<2>{
%     \item
%     Problèmes de bandits multi-joueurs décentralisés
%     \begin{itemize}
%       \item
%       \href{https://hal.inria.fr/hal-01629733/document}{\textcolor{darkgreen}{``Multi-Player Bandits
%       Revisited''}}, \textbf{L. Besson} \&
%       \href{http://chercheurs.lille.inria.fr/ekaufman/research.html}{E.
%       Kaufmann},
%       \href{http://www.cs.cornell.edu/conferences/alt2018/accepted.html}{ALT}
%       (\emph{Algorithmic Learning Theory}), avril 2018
%       % \href{https://hal.inria.fr/hal-01629733}{HAL-01629733}.
%     \end{itemize}
%     }

%     \vspace*{10pt}
%     \uncover<3>{
%     \item
%     Problèmes non stationnaires, stationnaires par morceaux
%     \begin{itemize}
%       \item
%       \href{https://hal.inria.fr/hal-02006471/document}{\textcolor{darkgreen}{Analyse non asymptotique d'un test séquentiel de détection de ruptures et application aux bandits non stationnaires}}\\
%       \textbf{L. Besson} \&
%       \href{http://chercheurs.lille.inria.fr/ekaufman/research.html}{E.
%       Kaufmann}, GRETSI, août 2019
%       % \href{https://hal.inria.fr/hal-02006471}{HAL-02006471}.
%     % \item
%     %   \href{https://hal.inria.fr/hal-02006471/document}{\textcolor{darkgreen}{``The Generalized Likelihood Ratio Test meets klUCB: an Improved Algorithm for Piece-Wise Non-Stationary Bandits''}}\\
%     %   \textbf{L. Besson} \&
%     %   \href{http://chercheurs.lille.inria.fr/ekaufman/research.html}{E.
%     %   Kaufmann},
%     %   % \href{https://hal.inria.fr/hal-02006471}{HAL-02006471}
%     %   \href{https://arxiv.org/abs/1902.01575}{\textcolor{blue}{arXiv:1902.01575}}
%     \end{itemize}
%     }
%   \end{enumerate}

% \end{frame}

% \setbeamercovered{invisibile}

% \begin{frame}{Aperçu de mes publications - FIXME}

% International Conferences
% \begin{itemize}
% \item
%   \href{https://hal.inria.fr/hal-XXX/document}{Decentralized Spectrum Learning for IoT Wireless Networks Collision Mitigation}, by \href{https://moychris.wordpress.com/}{C. Moy} \& \textbf{L. Besson},
%   \href{https://sites.google.com/view/ISIoT2019/}{ISIoT Workshop} at \href{https://dcoss.org/}{DCOSS}, Santorini, Greece, May 2019.
%   % \href{https://hal.inria.fr/hal-XXX}{HAL-XXX}.
% \item
%   \href{https://perso.crans.org/besson/articles/BMBBM__IEEE_WCNC__2019.pdf}{Upper-Confidence
%   Bound for Channel Selection in LPWA Networks with Retransmissions},
%   by \href{https://remibonnefoi.wordpress.com/}{R. Bonnefoi}, \textbf{L. Besson}, J.
%   Manco-Vasquez,
%   \href{https://moychris.wordpress.com/}{C. Moy}.
%   \href{https://sites.google.com/view/wcncworkshop-motion2019/}{MoTION Workshop} at
%   \href{http://wcnc2019.ieee-wcnc.org/}{IEEE WCNC},
%   Marrakech, Maroc, April 2019, \href{https://hal.inria.fr/hal-02049824}{HAL-02049824}.
% \item
%   \href{https://perso.crans.org/besson/articles/BBM__Demo_ICT_2018.pdf}{MALIN:
%   « Multi-Arm bandit Learning for Iot Networks » with GRC: A TestBed
%   Implementation and Demonstration that Learning Helps}, by \textbf{L.
%   Besson}, \href{https://remibonnefoi.wordpress.com/}{R. Bonnefoi},
%   \href{https://moychris.wordpress.com/}{C. Moy}. Demonstration
%   presented at ICT (\emph{International Conference on Communication}),
%   Saint-Malo, France, June 2018. Cf. \texttt{YouTu.be/HospLNQhcMk}, \href{https://hal.inria.fr/hal-02013866}{HAL-02013866}.
% \end{itemize}


% International Conferences

% \begin{itemize}
% \item
%   \href{https://perso.crans.org/besson/articles/BBM__IEEE_WCNC_2019.pdf}{GNU
%   Radio Implementation of MALIN: « Multi-Armed bandits Learning for
%   Internet-of-things Networks »}, by \textbf{L. Besson},
%   \href{https://remibonnefoi.wordpress.com/}{R. Bonnefoi},
%   \href{https://moychris.wordpress.com/}{C. Moy}.
%   \href{http://wcnc2019.ieee-wcnc.org/}{IEEE WCNC},
%   Marrakech, Maroc, April 2019, \href{https://hal.inria.fr/hal-02006825}{HAL-02006825}.
% \item
%   \href{https://hal.inria.fr/hal-01629733/document}{Multi-Player Bandits
%   Revisited}, by \textbf{L. Besson} \&
%   \href{http://chercheurs.lille.inria.fr/ekaufman/research.html}{E.
%   Kaufmann}.
%   \href{http://www.cs.cornell.edu/conferences/alt2018/accepted.html}{ALT}
%   (\emph{Algorithmic Learning Theory}), Lanzarote, Spain, April 2018,
%   \href{https://hal.inria.fr/hal-01629733}{HAL-01629733}.
% \item
%   \href{https://hal.inria.fr/hal-01705292/document}{Aggregation of
%   Multi-Armed Bandits learning algorithms for Opportunistic Spectrum
%   Access}, by \textbf{L. Besson},
%   \href{http://chercheurs.lille.inria.fr/ekaufman/research.html}{E.
%   Kaufmann} \& \href{https://moychris.wordpress.com/}{C. Moy}.
%   \href{http://wcnc2018.ieee-wcnc.org/}{IEEE WCNC}, Barcelone, Spain,
%   avril 2018, \href{https://hal.inria.fr/hal-01705292}{HAL-01705292}.
% \item
%   \href{https://hal.inria.fr/hal-01575419/document}{Multi-Armed Bandit
%   Learning in IoT Networks and non-stationary settings}, by
%   \href{https://remibonnefoi.wordpress.com/}{R. Bonnefoi}, \textbf{L.
%   Besson}, \href{https://moychris.wordpress.com/}{C. Moy},
%   \href{http://chercheurs.lille.inria.fr/ekaufman/research.html}{E.
%   Kaufmann} \& J. Palicot. \href{http://crowncom.org/2017/}{CrownCom}
%   (\emph{Conference on Cognitive Radio Oriented Wireless Networks}),
%   Lisbonne, Portugal, September 2017,
%   \href{https://hal.inria.fr/hal-01575419}{HAL-01575419}, \textbf{Best Paper Award}.
% \end{itemize}



% Other Publications

% \begin{itemize}
% \item
%   \href{https://hal.inria.fr/hal-02006471/document}{Analyse non asymptotique d'un test séquentiel de détection de ruptures et application aux bandits non stationnaires}, by \textbf{L. Besson} \&
%   \href{http://chercheurs.lille.inria.fr/ekaufman/research.html}{E.
%   Kaufmann}, March 2019,
%   \href{https://hal.inria.fr/hal-02006471}{HAL-02006471}.
% \item
%   \href{https://hal.inria.fr/hal-02006471/document}{The Generalized Likelihood Ratio Test meets klUCB: an Improved Algorithm for Piece-Wise Non-Stationary Bandits}, by \textbf{L. Besson} \&
%   \href{http://chercheurs.lille.inria.fr/ekaufman/research.html}{E.
%   Kaufmann}, February 2019,
%   \href{https://hal.inria.fr/hal-02006471}{HAL-02006471}.
% \item
% 	\href{https://hal.inria.fr/hal-01840022/document}{SMPyBandits: an
% 	Open-Source Research Framework for Single and Multi-Players Multi-Arms
% 	Bandits (MAB) Algorithms in Python}, by \textbf{L. Besson},
% 	active development since 2016,
% 	\href{https://hal.inria.fr/hal-01840022}{HAL-01840022}. Code on
% 	\href{https://GitHub.com/SMPyBandits/SMPyBandits}{GitHub.com/SMPyBandits}
% 	(about 40000 lines), documentation and tutorials on
%   \href{https://SMPyBandits.rtfd.io}{SMPyBandits.rtfd.io}.
% \item
% 	\href{https://hal.inria.fr/hal-01736357/document}{What Doubling-Trick
% 	Can and Can't Do for Multi-Armed Bandits}, by \textbf{L. Besson} \&
% 	\href{http://chercheurs.lille.inria.fr/ekaufman/research.html}{E.
% 	Kaufmann}, September 2018,
%   \href{https://hal.inria.fr/hal-01736357}{HAL-01736357}.

% \end{itemize}

% \end{frame}


% \subsection{\hfill{}Algorithme \emph{``Upper Confidence Bound''}\hfill{}}

% \begin{frame}{Exemple de solution : l'algorithme UCB}

% Au temps $t$, on joue le bras $C(t)\in\{1,\dots,K\}$, on observe la récompense \emph{iid} aléatoire $r(t) \sim \mathrm{Bernoulli}(\mu_k) \in\{0,1\}$.
% But : maximiser $\sum\limits_{t=1}^T r(t)$

% \pause
% \begin{itemize}
%   \item
%   Compter le nombre d'échantillons et la somme des récompenses pour chaque bras
%   % \item
%   % Et la somme des récompenses
%   \[
%     N_k(t) := \sum\limits_{s < t} \bold{1}(C(s)=k),
%     \;\;\;\;
%     X_k(t) := \sum\limits_{s < t} r(s) \bold{1}(C(s)=k)
%   \]
%   \pause
%     \item
%   Calculer $\mathrm{UCB}_k(t) = X_k(t) / N_k(t) + \sqrt{\alpha \log(t) / N_k(t)}$ \\
%   $=$ une borne supérieure de confiance sur la moyenne $\mu_k$ \alert{inconnue}
%   \item
%   Jouer le bras d'UCB maximum : $C(t) = \arg\max_k \mathrm{UCB}_k(t)$
%   % \\
%   % $\hookrightarrow$ Principe d'optimisme face à l'incertain
%     \item
%   % $\alpha$ contrôle l'
%   Équilibre entre \emph{exploitation} ($\alpha\to0$) et \emph{exploration}  ($\alpha\to\infty$)
% \end{itemize}

% \end{frame}


% \subsection{\hfill{}Mesure de performance pour les algorithmes de bandits\hfill{}}

\begin{frame}{Mesurer la performance d'un algorithme $\mathcal{A}$ avec son regret (moyen) $R_{\mathcal{A}}(T)$}

\pause

\begin{itemize}
  \item
  Différence de récompenses accumulées entre un ``oracle'' et $\mathcal{A}$

  \item
  L'algorithme ``oracle'' joue le meilleur bras $k^* = \arg\max \mu_k$
  %  pour tout $t$
  % $\Longrightarrow$
  % accumule les récompenses (moyennes) $\mathbb{E}\left[ \sum\limits_{t=1}^T r(t) \right] = T \mu_{k^*}$

  % \item
  % Algorithme $\mathcal{A}$ : accumule $\sum\limits_{t=1}^T \mathbb{E}\left[ r(t) \right]$

  \item
  Maximiser les récompenses cumulées
  $\Longleftrightarrow$ \alert{minimiser le regret}
  %
  \[ \alert{ R_{\mathcal{A}}(T) } := \mathbb{E}\left[ \sum\limits_{t=1}^T r_{k^*}(t) \right] - \sum\limits_{t=1}^T \mathbb{E}\left[ r(t) \right] = T \mu_{k^*} - \sum\limits_{t=1}^T \mathbb{E}\left[ r(t) \right]. \]

\end{itemize}

\pause
\vspace*{10pt}

\begin{exampleblock}{Régime typique pour des problèmes stationnaires (borne inf \& sup)}
  \begin{itemize}
  \item
  Aucun algorithme ne peut obtenir mieux que
  \hfill{}
  $R_{\mathcal{A}}(T) \geq \Omega(\log(T))$

  \item
  Et un algorithme efficace $\mathcal{A}$ obtient
  \hfill{}
  $R_{\mathcal{A}}(T) \leq \mathcal{O}(\log(T))$
  \end{itemize}
\end{exampleblock}

% \pause

% \begin{exampleblock}{Exemple de borne de regret : pour l'algorithme UCB}
%   Pour tout problème avec $K$ bras distribués selon des lois de Bernoulli, de paramètres $\mu_1,\dots,\mu_K \in[0,1]$, de bras optimal $k^*$
%   \begin{small}
%     \[ R_T^{\mathrm{UCB}} \leq \left( \sum_{\substack{k=1,\dots,K \\ \mu_k < \mu_{k^*}}} \frac{8}{(\mu_k - \mu_{k^*})^2} \right) \log(T) + o(\log(T)) = \mathcal{O}\left( \alert{C(\mu_1,\dots,\mu_K)} \log(T) \right). \]
%   \end{small}%
% \end{exampleblock}

\end{frame}


% \begin{frame}[plain]{Simulation de problèmes de bandits stationaires, avec ma bibliothèque SMPyBandits (en Python)}

%   % \centering
%   \includegraphics[width=1.20\textwidth]{figures/example_regret_stationary_bandits.pdf}

% \end{frame}


% \begin{frame}[plain]{Sélection dynamique du canal radio pour des objets communicants ``IoT'' (Internet of Things)}

%   \centering
%   \includegraphics[width=0.85\textwidth]{figures/Layout.pdf}

%   \vspace*{-15pt}
%   \textbf{Modèle} :
%   bras $=$ canal radio,
%   message reçu/perdu $=$ récompense $1$/$0$

% \end{frame}


% \begin{frame}[plain]{Interface utilisateur d'une démonstration sur du matériel réel présentée à la conférence ICT 2018}

%   \centering
%   \includegraphics[width=1.00\textwidth]{figures/UI.pdf}

% \end{frame}


% \begin{frame}[plain]{Comparaison entre un accès uniforme (naïf) et un accès intelligent (algorithme UCB)}

%   \centering
%   \includegraphics[width=0.50\textwidth]{figures/Object2_UniformAccess__UI.png}
%   ~
%   \includegraphics[width=0.47\textwidth]{figures/Object1_UCB__UI.png}

%   Conclusion : utiliser un algorithme peu coûteux de ``bandits multi-bras'' améliore les performances de l'objet connecté (ici $40\% \Rightarrow 60\%$)

% \end{frame}


% \begin{frame}{Contributions principales de ma thèse}
\begin{frame}{Deux contributions principales de ma thèse}

Pour différentes extensions du modèle classique de bandits\ldots{}

\begin{itemize}
  \item
  Formalisation mathématique
  \item
  Nous proposons de nouveaux algorithmes\ldots{}
  % \item
  \newline
  avec une nouvelle analyse théorique, nouvelles bornes de regret\ldots{}
  % \item
  \newline
  validées par des simulations numériques\ldots{}
  \item
  $\implies$ Améliore l'état de l'art sur les deux aspects !
\end{itemize}

\end{frame}


\begin{frame}{Deux contributions principales de ma thèse}

\begin{exampleblock}{Bandits stationnaires (classiques) avec $K$ bras et $T$ étapes}
  \[ R_{\mathcal{A}}(T) = \mathcal{O}\left( K \log(T) \right). \]
\end{exampleblock}

% Pour différentes extensions du modèle classique de bandits\ldots{}

\begin{block}<2->{1) Bandits \textbf{multi-joueurs décentralisé} \hspace*{40pt} \href{https://arxiv.org/abs/1711.02317}{[Besson et al, ALT, 2018]}}
  Si \alert{$M \leq K$ joueurs} jouent face au même problème de bandit, \alert{avec collisions} mais sans communication entre eux ni sans contrôle centralisé :
  %
  $R_{\mathrm{MCTopM}}(T) = \mathcal{O}\left( K \alert{M^3} \log(T) \right)$.
\end{block}

% \pause

% \begin{block}{Analyse générique des \textbf<2>{``doubling trick''} \hspace*{55pt} \href{https://arxiv.org/abs/1803.06971}{[arXiv:1803.06971]}}
%   % % Analyse unifiée des ``doubling trick'' exponentiels et géométrique.
%   % But : si l'algorithme $\mathcal{A}$ doit connaître l'horizon $A$, une technique ``doubling trick'' le rend ``anytime'' (e.g., $T=2,4,8,\dots,2^i,\dots$),
%   % sans trop réduire ses garanties en regret :
%   Une \alert<2>{technique ``\emph{doubling trick}''} enlève la dépendance en $T$ d'un algorithme $\mathcal{A}$ non-\emph{anytime}, et garde l'ordre de grandeur de son regret :
%   %
%   \begin{small}
%     \[
%       % \forall 0<\gamma<1, \delta>0,
%       % \;\;
%       R_T^{\mathcal{A}_T} = \mathcal{O}\left( C_{\mathcal{A}} T^{\gamma} (\log(T))^{\delta} \right)
%       \implies
%       R_T^{\alert<2>{DT(\mathcal{A})}} = \mathcal{O}\left( \alert<2>{\ell(\gamma,\delta) \times} C_{\mathcal{A}} T^{\gamma} (\log(T))^{\delta} \right).
%     \]
%   \end{small}%
% \end{block}

% \pause

\begin{block}<3->{2) Bandits stationnaires \textbf{par morceaux} \hspace*{25pt} \href{https://arxiv.org/abs/1902.01575}{[Besson et al, GRETSI, 2019]}}
  Si le problème est \alert{stationnaire par morceaux, sur $\Upsilon = o(\sqrt{T})$ intervalles ``assez grands''} :
  %
  $R_{\mathrm{B}\text{-}\mathrm{GLR}}(T) = \mathcal{O}\left( K \sqrt{\alert{\Upsilon} T \log(T)} \right)$.
\end{block}

\end{frame}


% \begin{frame}[plain]{Bandits \alert{multi-joueurs} [Besson et al, ALT, 2018]}

%   % \centering
%   \includegraphics[width=1.20\textwidth]{figures/example_regret_multi_players_bandits.png}

% \end{frame}


\begin{frame}[plain]{\normalsize Simulations d'algorithmes de bandits \alert{non stationnaires} avec ma bibliothèque SMPyBandits, écrite en Python \hfill{}\textcolor{gray}{[Besson et al, GRETSI, 2019]}}

  \centering
  \includegraphics[width=1.15\textwidth]{figures/example_regret_non_stationary_bandits.png}

\end{frame}


% \subsection{\hfill{}Thèse bientôt terminée\hfill{}}

\begin{frame}{Calendrier de la fin de ma thèse}

% \begin{block}{Je termine ma thèse de doctorat}

\begin{itemize}
% \tightlist
\item
  Rédaction bientôt terminée
  \newline
  $\implies$ \textbf{Manuscrit envoyé en juillet}
% \item
%   Mon manuscrit est basé sur mes principales publications
\vspace*{20pt}
\item
  Jury constitué
  \newline
  Rapporteurs en attente de ma thèse
  \newline
  $\implies$ \textbf{Soutenance prévue en octobre}
\end{itemize}

% \end{block}

% \begin{block}{D'ici septembre 2019}

% \begin{itemize}
% \tightlist
% \item
%   Un séjour à Lille fin juin pour assister aux oraux d'agrégation ?\\
%   Si j'ai le temps\ldots{} \\
%   \(\Longrightarrow\) important de garder une expérience actuelle du concours !
% \item
%   Je devrais présenter un article au GRETSI 2019 (Lille, fin août)
% \end{itemize}

% \end{block}

\end{frame}



% \begin{frame}[plain]{Aperçu de mon manuscrit
% (\href{https://GitHub.com/Naereen/phd}{GitHub.com/Naereen/phd})}

% % \begin{figure}
% % \centering
% \includegraphics{figures/apercu_plan_these.png}
% % \caption{bg original 90\%}
% % \end{figure}

% \end{frame}



\subsection{\hfill{}3.3. Mon projet de recherche à l'IRISA\hfill{}}

\begin{frame}{Mon projet de recherche \only<2->{à l'IRISA}}

% But : poursuivre une \textbf{activité de recherche} après ma thèse, \textbf{en informatique et mathématiques}, en ``machine learning'' et algorithmique numérique,
% mais en changeant de domaine d'applications

But : continuer à résoudre \alert{\textbf{des problèmes de détection et d'adaptation automatique à la non stationnarité}}

Et les appliquer à d'autres flux de données réels (ex : sons, images)\ldots

% \begin{itemize}
% \item
%   Après ma thèse, je souhaite poursuivre une activité de recherche
% \item
%   Toujours
%   % , mais sans contrainte quant aux applications
% \end{itemize}

\vspace*{15pt}
\pause

$\implies$ \textbf{L'équipe PANAMA est prête à m'accueillir à l'IRISA} :
\begin{itemize}
\item
  Rémi Gribonval et Frédéric Bimbot soutiennent mon projet
  % intégration dans leur équipe
\pause
\item
  Je vais \textbf{collaborer avec Jérémy Cohen},
  un chargé de recherche (CNRS) recruté en octobre 2018,
  spécialiste des méthodes de décomposition et factorisation matricielles et tensorielles
% \item
%   Nous avons identifié trois axes de collaboration
\end{itemize}

% \begin{block}{Organisation / prévision}

% \begin{itemize}
% \tightlist
% \item
%   1/2 jours par semaines à plein temps dès octobre
% \item
%   Une place dans les bureaux de l'équipe
% \item
%   Lectures et discussions très régulières, octobre-décembre
% \item
%   Recherche active dès janvier 2020
% \item
%   Je pense pouvoir être ``productif'' dès mars prochain
% \end{itemize}

% \end{block}

\end{frame}


\begin{frame}{3 axes de collaboration avec Jérémy Cohen}

\begin{enumerate}[1.]
\item
  \alert<1>{Méthodes matricielles et tensorielles étendues au \textbf<1>{cadre séquentiel}} \\
  \(\hookrightarrow\) mélange nos deux domaines d'expertise
\vspace*{20pt}
\pause
\item
  \alert<2>{Détection et adaptation automatique à des \textbf<2>{non-stationnarités}} \\
  \(\hookrightarrow\) encore trop peu étudié pour ce genre de problèmes, mais très utile ``en pratique''
\vspace*{20pt}
\pause
\item
  \alert<3>{Sélection automatique du meilleur algorithme et \textbf<3>{hyper apprentissage} automatique (séquentiel ou non), avec des bandits} \\
  \(\hookrightarrow\) déjà très étudié, mais pas encore dans ce domaine
  % appliqué aux méthodes matricielles/tensorielles
\end{enumerate}

% \pause

% Suite logique de deux des questions au cœur de ma thèse :
% \begin{itemize}\tightlist
%   \item
%   agrégation séquentielle d'algorithmes de bandits \hfill{}\textcolor{gray}{[Besson \& Moy \& Kaufmann, WCNC 2018]}
%   \item
%   test de vraisemblance généralisé pour un algorithme adaptatif qui résout les problèmes de bandits stationnaires par morceaux
%   % \hfill{}\textcolor{gray}{[arXiv:1902.01575]}
% \end{itemize}

\end{frame}



% \begin{frame}{Autres aspects de mon intégration à l'IRISA}

% \begin{enumerate}
% [1.]
% \tightlist
% \item
%   Séminaire de l'équipe PANAMA

%   \begin{itemize}
%   \tightlist
%   \item
%     J'y présente mon dernier travail début juin !
%   \end{itemize}
% \item
%   Créer un groupe de lecture ``Deep Reinforcement Learning''

%   \begin{itemize}
%   \tightlist
%   \item
%     Possiblement lié au cours d'apprentissage statistique à l'ENS ?
%   \end{itemize}
% \item
%   Formations techniques : deux demi-journées annuelles (en 2020)

%   \begin{itemize}
%   \tightlist
%   \item
%     diffusion et production de codes libres,
%   \item
%     tests et documentations de codes de recherche en mathématiques
%     appliquées et informatique
%   \end{itemize}
% \item
%   Encadrer un-e stagiaire de recherche à partir d'avril 2020 ?
% \end{enumerate}

% \end{frame}

% \begin{frame}{%
% \protect\hypertarget{autres-objectifs-personnels}{%
% 5. Autres objectifs personnels}}

% \begin{enumerate}
% [1.]
% \tightlist
% \item
%   Terminer mes travaux de thèse en cours (un ou deux articles à réécrire
%   et soumettre ailleurs ?)
% \end{enumerate}

% \begin{enumerate}
% [1.]
% \setcounter{enumi}{1}
% \tightlist
% \item
%   Remettre un pied en \textbf{classe préparatoire} ?

%   \begin{itemize}
%   \tightlist
%   \item
%     En septembre, peut-être demander une autorisation de cumul ?
%   \item
%     pour des khôlles de maths et/ou TD/TP d'informatique
%   \item
%     \emph{maximum 2h/semaine}, si la charge d'AGPR le permet
%   \item
%     \(\hookrightarrow\) en MPSI/MP, avec Amélie Stainer au Lycée
%     Chateaubriand ?
%   \end{itemize}
% \end{enumerate}

% \begin{enumerate}
% [1.]
% \setcounter{enumi}{2}
% \tightlist
% \item
%   Rester à Rennes 2 ou 3 ans, voire plus !
% \end{enumerate}

% \end{frame}

% \begin{frame}{%
% \protect\hypertarget{perspectives-apruxe8s-ce-poste}{%
% Perspectives (après ce poste)}}

% Après 2/3\(^+\) années en tant qu'AGPR à l'ENS de Rennes\ldots{}

% \begin{block}{Prochain poste ?}

% \begin{itemize}
% \item
%   Mouvement spécifique (EN), pour demander un poste en classe
%   "prépa"ratoire (maths+info / info)
% \item
%   Peut-être aussi postuler au CNRS, Inria ou Maître de Conférence
%   (sections 26, 27) ?
% \end{itemize}

% \end{block}

% \begin{block}{\emph{Mais aussi ?}}

% \begin{itemize}
% \tightlist
% \item
%   Enseigner l'informatique dans le secondaire\ldots{}?
% \item
%   Lycée Français à l'étranger\ldots{}?
% \item
%   Retourner enseigner à l'étranger (écoles d'ingénieurs, prépa
%   etc)\ldots{}?
% \end{itemize}

% \end{block}

% \end{frame}

\section{\hfill{}Conclusion\hfill{}}
\subsection{}


\begin{frame}{Conclusion}

\begin{center}
  % \raggedright
  \includegraphics[width=0.80\linewidth,right]{figures/merci_ecriture_manuelle.png}
\end{center}

% \begin{center}
%   \begin{Large}
%     Seconde partie\\
%     Questions \& discussion
%   % {\Fontify Merci de votre attention .}
%   % \Smiley[0.9]
%   \end{Large}
% \end{center}

\end{frame}

\end{document}
